% ============================================================
%  Symbolic-JEPA: Domain-Adaptive World Models for Physics Reasoning
%  arXiv preprint — Vortaz Labs, April 2026
% ============================================================
\documentclass[10pt,twocolumn]{article}

\usepackage[margin=0.85in,top=1in,columnsep=0.25in]{geometry}
\usepackage{amsmath,amssymb,amsthm}
\usepackage{graphicx}
\usepackage{booktabs}
\usepackage{multirow}
\usepackage{hyperref}
\usepackage[T1]{fontenc}
\usepackage{microtype}
\usepackage{xcolor}
\usepackage{enumitem}
\usepackage{caption}
\usepackage{subcaption}
\usepackage{url}
\usepackage{times}
\usepackage{natbib}
\usepackage{mdframed}

\definecolor{myblue}{RGB}{29,78,216}
\definecolor{myred}{RGB}{185,28,28}
\definecolor{mygreen}{RGB}{6,95,70}
\definecolor{mygray}{RGB}{107,114,128}

\hypersetup{colorlinks,linkcolor=myblue,citecolor=myblue,urlcolor=myblue}
\setlength{\abovecaptionskip}{4pt}
\setlength{\belowcaptionskip}{2pt}

% ── Inline concept box ────────────────────────────────────────────────────────
\newmdenv[
  backgroundcolor=gray!8,
  linecolor=gray!40,
  linewidth=0.6pt,
  innertopmargin=4pt,
  innerbottommargin=4pt,
  innerleftmargin=6pt,
  innerrightmargin=6pt,
  skipabove=4pt,
  skipbelow=4pt,
]{conceptbox}

\title{\textbf{Symbolic-JEPA: Domain-Adaptive World Models\\for Physics Reasoning}}
\author{
  \textbf{Ishaan Verma}\\
  Vortaz Labs\\
  \texttt{ishaan@vortazlabs.ai}
}
\date{April 2026}

\begin{document}
\maketitle

% ── Abstract ─────────────────────────────────────────────────────────────────
\begin{abstract}
We introduce \textbf{Symbolic-JEPA (S-JEPA)}, a domain-adaptive world model that extends
the Joint-Embedding Predictive Architecture~\citep{lecun2022path,assran2023ijepa} to
structured physical symbols rather than pixels. A single ${\sim}2$M-parameter S-JEPA,
trained in 26~minutes on a consumer GPU, achieves $1.8\times$ the score of a
1B-parameter LLM fine-tuned on \emph{identical} trajectory data in live Unity physics gameplay,
and $16\times$ the score of the best zero-shot 120B-parameter LLM.
The fine-tuned LLM exhibits \textbf{Behavioral Collapse} — degenerating to a fixed action
regardless of scene state — a failure that persists across model scales from 1B to 120B
parameters in this setting, consistent with a structural mismatch between token prediction
and continuous spatial dynamics rather than a data or scale limitation.
The same frozen S-JEPA core transfers to scientific CFD aerodynamics prediction
(PLAID AirfRANS, zero-shot NACA-5 transfer at $0.986$ cosine similarity)
and 2D planar manipulation (28\% fewer trainable parameters, negligible performance loss).
\end{abstract}

% ── 1. Introduction ───────────────────────────────────────────────────────────
\section{Introduction}
\label{sec:intro}

Physical reasoning — predicting how objects move, collide, and interact under applied forces
— is fundamental to robotics, scientific simulation, and embodied AI. A system that
\emph{genuinely understands physics} should interpolate that a ball at $9.3^\circ$ follows a
different trajectory than one at $9.4^\circ$; no discrete token captures this.

Large language models (LLMs)~\citep{brown2020gpt3,touvron2023llama2} are optimised for
predicting the next \emph{token} — a process suited to discrete symbolic reasoning but
structurally mismatched to continuous physical dynamics. A dominant counter-argument holds
that fine-tuning resolves this. We test it directly: we fine-tune Llama-3.2-1B via Unsloth
4-bit QLoRA~\citep{dettmers2023qlora} on the \emph{same} 50K physics trajectories as
S-JEPA. The result is \textbf{Behavioral Collapse} — the model fires at $0.2^\circ$/100\%
power on every shot regardless of scene state.

S-JEPA follows LeCun's world model vision~\citep{lecun2022path}: predict in a compressed
\emph{latent space}, then select actions by gradient-based optimisation through this internal
model. Building on I-JEPA~\citep{assran2023ijepa} and V-JEPA~\citep{bardes2024vjepa},
S-JEPA replaces the pixel encoder with a domain-specific \emph{symbol encoder} operating on
compact vectors of physical attributes.

\paragraph{Contributions.}
\begin{enumerate}[leftmargin=*,topsep=2pt,itemsep=1pt]
\item \textbf{S-JEPA}: a symbolic JEPA with an energy-based critic and multi-start gradient
      actor, validated on three physically heterogeneous domains in 26~min of training.
\item \textbf{Behavioral Collapse}: characterisation of a systematic failure mode in
      fine-tuned LLMs under continuous physics execution — observed across model scales
      1B to 120B — with matched-data evidence consistent with an architectural explanation.
\item \textbf{Cross-domain transfer}: frozen Predictor + Critic transfers from game physics to
      CFD and planar manipulation with 28\% fewer trainable parameters and negligible loss.
\end{enumerate}

% ── 2. Background ─────────────────────────────────────────────────────────────
\section{Background and Related Work}
\label{sec:background}

\paragraph{JEPA.}
JEPA~\citep{lecun2022path} trains world models to predict in \emph{latent space} rather than
pixel space, avoiding capacity waste on irrelevant details. I-JEPA~\citep{assran2023ijepa}
applies this to static images; V-JEPA~\citep{bardes2024vjepa} to video. Both use pixel
encoders. S-JEPA instead processes compact \emph{symbol vectors} of physical attributes —
positions, velocities, forces — directly.

\begin{conceptbox}
\textbf{Why latent-space prediction?} Predicting raw future pixels requires reconstructing
every detail (background, lighting) regardless of physical relevance. Predicting a 256D
latent vector focuses capacity entirely on physical structure.
\end{conceptbox}

\paragraph{Energy-Based Models.}
An EBM~\citep{lecun2006ebm} assigns a scalar energy to input pairs: low energy = compatible,
high energy = incompatible. S-JEPA's Energy Critic scores predicted next states; the Actor
minimises this energy to find good actions.

\paragraph{World Models for Control.}
Ha \& Schmidhuber~\citep{ha2018worldmodels} and DreamerV3~\citep{hafner2023dreamerv3}
plan in pixel-space latent representations. S-JEPA operates on symbolic inputs and uses
gradient-based action search (not RL), making it sample-efficient and transferable.

\paragraph{LLMs for Physics.}
LLMs have been applied to game-playing~\citep{brown2020gpt3} and robot
planning~\citep{ahn2022saycan}, but prior work suggests they struggle with precise spatial
dynamics. We provide systematic fine-tuned evidence under the strongest fair comparison.

\paragraph{Scientific Surrogates.}
PINNs~\citep{raissi2019pinn} and neural operators~\citep{li2020fno} approximate physical
PDEs. AirfRANS~\citep{bonnet2022airfrans} provides standardised CFD evaluation against
OpenFOAM.

\paragraph{Physics Neural Simulators.}
Graph Network Simulators~\citep{sanchez2020gns} and MeshGraphNets~\citep{pfaff2021meshgraphnets}
learn particle or mesh dynamics via message-passing GNNs, achieving high fidelity on
structured physical systems. Both require per-instance graph construction and full
retraining when the domain or resolution changes.
S-JEPA's fixed-size symbolic encoding trades some in-distribution fidelity for
parameter-efficient zero-shot transfer and gradient-based action search — a deliberate
design point for multi-domain deployment from limited data.

\paragraph{Structured World Models.}
Contrastive Structured World Models~\citep{kipf2020cswm} learn object-centric latent
representations via contrastive prediction. S-JEPA extends this direction with an
energy-based critic, EMA target encoder, and SIGReg variance regularisation, and is
the first to validate such a model across three physically heterogeneous domains.

% ── 3. Architecture ───────────────────────────────────────────────────────────
\section{Symbolic-JEPA Architecture}
\label{sec:arch}

Figure~\ref{fig:arch} shows the complete S-JEPA architecture: Encoder, Target Encoder
(EMA copy), Predictor, Energy Critic, and Decoder. At inference, only Encoder, Predictor,
and Critic are active.

\begin{figure*}[t]
\centering
\includegraphics[width=\linewidth]{sjepa_architecture.pdf}
\caption{
\textbf{S-JEPA Architecture.}
\textit{Inference path (top):} The Encoder maps domain symbols to a 256D latent $z_t$.
The Predictor combines $z_t$ with action $a_t$ to predict $\hat{z}_{t+1}$.
The Energy Critic scores the predicted state; the Actor searches for the action minimising
this energy via multi-start gradient descent.
\textit{Training only (bottom):} A momentum-updated Target Encoder produces stable
regression targets $z_{t+1}$; JEPA loss minimises $\|\hat{z}_{t+1} - z_{t+1}\|^2$.
The Decoder reconstructs symbols from $z_t$ as an auxiliary task.
\textbf{The Encoder (amber) is the only component swapped between domains.}
}
\label{fig:arch}
\end{figure*}

\subsection{Symbolic State Representation}

Rather than processing pixels through a CNN or ViT, S-JEPA takes a compact \emph{symbol
vector}: a fixed-size array encoding object positions, types, and physical attributes.

\begin{conceptbox}
\textbf{Why symbols instead of pixels?}
A pixel image of an Angry Birds scene contains ${\sim}10^6$ numbers; the physically relevant
information (pig positions, block geometry) fits in ${\sim}100$ numbers — $10{,}000\times$
more compact, making the latent space far easier to learn.
\end{conceptbox}

The three domains use different symbol vectors:
\begin{itemize}[leftmargin=*,topsep=2pt,itemsep=1pt]
\item \textbf{Science Birds} (2D rigid-body, Unity): $\mathbf{s} \in \mathbb{R}^{164}$
  encoding $(x,y)$ coordinates, object type, bounding box, and connectivity for up to
  20 scene entities, extracted via PIL connected-component detection.
\item \textbf{PLAID AirfRANS} (CFD): per-mesh-point features $(u,v,d,n_x,n_y)$ aggregated
  to a 5D global feature vector capturing mean flow and boundary geometry.
\item \textbf{Planar Grasping} (PyBullet): $\mathbf{s} \in \mathbb{R}^{32}$ encoding
  end-effector pose, joint angles, fingertip forces, object offset, and contact geometry.
\end{itemize}

\subsection{Encoder and SIGReg}

The Encoder $f_\phi : \mathbb{R}^d \to \mathbb{R}^{256}$ is a two-hidden-layer MLP:
\begin{equation}
z_t = \mathrm{LN}\!\left(W_2\,\sigma\!\left(W_1\,\mathbf{s}_t + b_1\right) + b_2\right)
\label{eq:encoder}
\end{equation}
where $\sigma$ is GELU activation and $\mathrm{LN}$ is Layer Normalisation (stabilises
training by normalising output to zero mean, unit variance).

\textbf{The Encoder is the only domain-specific component.} All other parameters are frozen
at transfer.

Self-supervised training is prone to \emph{dimensional collapse}: all inputs map to the
same latent point. We prevent this with SIGReg variance regularisation:
\begin{equation}
\mathcal{L}_{\mathrm{reg}} = \sum_{j=1}^{256} \max\!\left(0,\, \gamma - \mathrm{Std}(z^j_{\mathrm{batch}})\right)
\label{eq:sigreg}
\end{equation}
Any latent dimension whose standard deviation falls below $\gamma = 1$ is penalised,
encouraging every dimension to carry useful information.
SIGReg is related to the variance term of VICReg~\citep{bardes2022vicreg} but applies
per-dimension standard deviation rather than full batch covariance, reducing compute
from $O(d^2)$ to $O(d)$ and remaining numerically stable at batch size 256.

\subsection{EMA Target Encoder}

A second copy of the Encoder provides stable prediction targets via Exponential Moving
Average (EMA) — a slow momentum update that prevents training instability:
\begin{equation}
\bar{\phi} \leftarrow \mu\,\bar{\phi} + (1-\mu)\,\phi, \qquad \mu = 0.98
\label{eq:ema}
\end{equation}
The Target Encoder receives \emph{no gradient} and is discarded at inference.

\subsection{Predictor}

The Predictor $g_\psi : \mathbb{R}^{256+k} \to \mathbb{R}^{256}$ is the world model
core — it predicts the next latent state given the current state and action:
\begin{equation}
\hat{z}_{t+1} = g_\psi\!\left([z_t \;\|\; a_t]\right)
\label{eq:predictor}
\end{equation}
The action vector $a_t$ encodes: $(\alpha, p, \tau)$ for Science Birds (angle, power,
tap time; $k{=}3$); gripper pose for grasping ($k{=}3$); inlet velocity and angle
of attack for CFD ($k{=}5$). The Predictor is a two-layer MLP
($256{+}k \to 512 \to 256$) with GELU and Layer Normalisation.

\subsection{Energy Critic}

The Critic $E_\omega : \mathbb{R}^{256} \to \mathbb{R}$ assigns a scalar energy to the
predicted next state — low energy = good outcome, high energy = bad:
\begin{equation}
e = E_\omega(\hat{z}_{t+1}) = \mathbf{w}^\top\,\sigma\!\left(W_3\,\sigma\!\left(W_2\,
\sigma\!\left(W_1\,\hat{z}_{t+1}\right)\right)\right)
\label{eq:critic}
\end{equation}
The Critic ($256 \to 512 \to 256 \to 1$) correlates strongly with game outcomes
(Spearman's $\rho = 0.71$, $p < 10^{-10}$) and OpenFOAM CFD ground truth after encoder
adaptation ($\rho = 0.509$, $p = 8.9\times 10^{-28}$).

\begin{conceptbox}
\textbf{Spearman's $\rho = 0.71$} means shots ranked higher by the Critic tend to score
higher in the game in 71\% of rank-order comparisons. $p < 10^{-10}$ means this is
extremely unlikely to arise by chance.
\end{conceptbox}

\subsection{Actor and Decoder}

The Actor finds the best action by gradient descent through the frozen Predictor and Critic:
\begin{equation}
a^* = \arg\min_{a \in \mathcal{A}} \; E_\omega\!\left(g_\psi\!\left([z_t \;\|\; a]\right)\right)
\label{eq:actor}
\end{equation}

\begin{conceptbox}
\textbf{Why multi-start?} The energy landscape over actions may have multiple local minima.
We run $N{=}8$ random starting points for 100 Adam steps each and pick the lowest-energy
result, reducing the chance of converging to a suboptimal action.
\end{conceptbox}

The auxiliary \textbf{Decoder} $h_\xi : \mathbb{R}^{256} \to \mathbb{R}^d$ reconstructs
the symbol vector from $z_t$, providing an additional training signal and enabling
interpretability analysis. It is dropped entirely at inference.

\subsection{Parameter Summary}

\begin{table}[h]
\centering
\small
\setlength{\tabcolsep}{4pt}
\begin{tabular}{lcr}
\toprule
Component & Architecture & Params \\
\midrule
Encoder (online) & $d \to 512 \to 256$ & $\sim$0.49M \\
Target Encoder   & EMA copy (no grad) & $\sim$0.49M \\
Predictor        & $259 \to 512 \to 256$ & $\sim$0.53M \\
Energy Critic    & $256 \to 512 \to 256 \to 1$ & $\sim$0.50M \\
Decoder          & $256 \to 512 \to d$ & $\sim$0.19M \\
\midrule
\textbf{Total (all)} & & \textbf{$\sim$2.17M} \\
\textit{Inference only} & \textit{Enc + Pred + Critic} & \textit{$\sim$1.52M} \\
\bottomrule
\end{tabular}
\caption{S-JEPA parameter counts. Target Encoder and Decoder are training-only.}
\label{tab:params}
\end{table}

% ── 4. Training Procedure ─────────────────────────────────────────────────────
\section{Training Procedure}
\label{sec:training}

\paragraph{JEPA Prediction Loss.}
Given a transition $(\mathbf{s}_t, a_t, \mathbf{s}_{t+1})$, train the Predictor to match
the Target Encoder's encoding of the true next state:
\begin{equation}
\mathcal{L}_{\mathrm{JEPA}} = \left\|\hat{z}_{t+1} - \mathrm{sg}(z_{t+1})\right\|_2^2
\label{eq:jepa_loss}
\end{equation}
$\mathrm{sg}(\cdot)$ is \emph{stop-gradient}: no backprop through the Target Encoder
(updated by EMA only).

\paragraph{Critic Contrastive Margin Loss.}
Train the Critic to assign lower energy to real next states $\hat{z}^+$ than to
counterfactuals $\hat{z}^{-,j}$ sampled from other batch trajectories:
\begin{equation}
\mathcal{L}_{\mathrm{critic}} = \frac{1}{K}\sum_{j=1}^{K}
\max\!\left(0,\; E_\omega\!\left(\hat{z}^+\right) - E_\omega\!\left(\hat{z}^{-,j}\right) + \Delta\right)
\label{eq:critic_loss}
\end{equation}
The hinge function ($\max(0,\cdot)$) with margin $\Delta{=}1.0$ is zero once the real
next state is already lower-energy than the counterfactual by a sufficient margin.
We use $2.3\times$ positive-sample oversampling for class balance.

\paragraph{Reconstruction \& Total Loss.}
The Decoder adds an $\ell_2$ reconstruction signal to anchor the latent space:
\begin{equation}
\mathcal{L} = \mathcal{L}_{\mathrm{JEPA}} + \lambda_c\,\mathcal{L}_{\mathrm{critic}}
              + \lambda_r\!\left\|h_\xi(z_t) - \mathbf{s}_t\right\|_2^2
              + \lambda_{\mathrm{reg}}\,\mathcal{L}_{\mathrm{reg}}
\label{eq:total_loss}
\end{equation}

\begin{table}[h]
\centering
\small
\setlength{\tabcolsep}{5pt}
\begin{tabular}{ll}
\toprule
Hyperparameter & Value \\
\midrule
Optimiser & Adam ($\beta_1{=}0.9$, $\beta_2{=}0.999$) \\
Learning rate & $3 \times 10^{-4}$ \\
Batch size & 256 transitions \\
Loss weights & $\lambda_c{=}0.5$, $\lambda_r{=}0.1$, $\lambda_\text{reg}{=}0.01$ \\
EMA momentum & $\mu = 0.98$ \\
Margin $\Delta$ & $1.0$ \\
Actor starts / steps & 8 / 100 \\
Training data & 50K transitions \\
Wall-clock time & 26~min (RTX 3090) \\
\bottomrule
\end{tabular}
\caption{Training hyperparameters for all Science Birds experiments.}
\label{tab:hparams}
\end{table}

% ── 5. Experimental Setup ─────────────────────────────────────────────────────
\section{Experimental Setup}
\label{sec:setup}

The same core architecture is used across all three domains; only the Encoder changes.

\subsection{Domain 1: Science Birds}

Science Birds~\citep{renz2019sciencebirds} is a Unity re-implementation of Angry Birds
for AI research, with a WebSocket interface. Pigs (targets) are placed among block
structures; a bird launched from a slingshot must destroy them. PhysX handles rigid-body
physics. We evaluate on 4 held-out levels of increasing difficulty (none appear in the 50K training rollouts).

State extraction uses a PIL-based connected-component detector producing the 164D symbol
vector. Action space: $a_t = (\alpha, p, \tau)$ — launch angle, power, and tap time —
all continuous. Data: 50K transitions collected via guided play and random exploration.

\subsection{Domain 2: PLAID AirfRANS (CFD)}

AirfRANS~\citep{bonnet2022airfrans} provides 200 OpenFOAM RANS simulations of 2D
NACA-4 airfoils. Each yields per-mesh-point velocity, pressure, and turbulent viscosity
fields. Given 5D input features $(u_\infty, v_\infty, d, n_x, n_y)$ (far-field velocity,
wall distance, surface normals), predict 4D field outputs. We also evaluate zero-shot on
NACA-5 airfoils (5-parameter family, never seen during training).

\subsection{Domain 3: Planar Manipulation (PyBullet)}

A synthetic 2D planar grasping dataset from PyBullet~\citep{coumans2016pybullet}.
All S-JEPA weights from Science Birds are frozen; only a new $32\text{D}{\to}256\text{D}$
Encoder is trained. This is a controlled transfer probe, not a real-robot evaluation.
Metric: cosine similarity between predicted and actual next latent states.

\subsection{Baselines}

\paragraph{Zero-shot LLMs.}
GPT-OSS-120B, GPT-OSS-20B, and Llama-3.1-8B receive structured text prompts describing
the scene and must output $(\alpha, p, \tau)$ as numbers.

\paragraph{Fine-tuned LLM (fairness protocol).}
Llama-3.2-1B fine-tuned via \textbf{Unsloth 4-bit QLoRA}~\citep{dettmers2023qlora}
on the \emph{same} 50K trajectories as S-JEPA, formatted as instruction-response JSONL pairs.
\emph{QLoRA} compresses weights to 4-bit integers then adds small trainable adapter matrices.

\paragraph{Classical ML (CFD only).}
Ridge Regression, Random Forest, and a 3-layer Deep MLP (256 units/layer) on the same
AirfRANS training split.

\paragraph{LLM regression format.}
Zero-shot and fine-tuned LLMs receive structured text prompts describing the scene state
and are required to output explicit floating-point action values $(\alpha, p, \tau)$.
Behavioral Collapse — outputting a fixed token sequence — is confirmed by identical action
outputs across all 4 scene configurations and all 3 evaluation runs, ruling out a
prompt-formatting artifact.

\paragraph{Reproducibility.}
Science Birds scores are averaged over $n{=}3$ independent random seeds (0, 1, 2).
CFD uses the standard AirfRANS 80/20 split (168 training / 32 test simulations).
Grasping uses 5{,}000 synthetic PyBullet episodes.
Inference latency per step: $<15$\,ms for Encoder\,+\,Predictor\,+\,Critic;
${\sim}800$\,ms for the 8-start gradient actor ($8 \times 100$ Adam steps at
${\approx}1$\,ms/step).
The 26-min training wall-clock excludes the offline data-collection phase.
Code and weights will be released at \url{https://github.com/vortazlabs/s-jepa}.

% ── 6. Results ────────────────────────────────────────────────────────────────
\section{Results}
\label{sec:results}

\subsection{Science Birds: Live Unity Leaderboard}

\begin{table}[h]
\centering
\small
\setlength{\tabcolsep}{3pt}
\begin{tabular}{lcccc}
\toprule
System & Params & Score ($\mu \pm \sigma$) & Levels \\
\midrule
\textbf{S-JEPA (ours)} & \textbf{$\sim$2M} & \textbf{2400 $\pm$ 80} & \textbf{3/4} \\
Llama-3.2-1B (FT) & 1B & 1325 $\pm$ 110 & $3/4^{\dagger}$ \\
GPT-OSS-120B (zero-shot) & 120B & 150 & 1/4 \\
GPT-OSS-20B (zero-shot) & 20B & 125 & 0/4 \\
Llama-3.1-8B (zero-shot) & 8B & 0 & 0/4 \\
\bottomrule
\end{tabular}
\caption{Live Unity leaderboard. $\dagger$: fine-tuned Llama cleared 3 levels via
collateral damage at fixed $0.2^\circ$/100\% power, not spatial targeting.}
\label{tab:sb}
\end{table}

\begin{figure}[t]
\centering
\includegraphics[width=\linewidth]{paper_figures/fig2_live_unity_scores.png}
\caption{
\textbf{Live Unity scores across 4 levels.}
S-JEPA ($\sim$2M params) scores $1.8\times$ the fine-tuned Llama-1B trained on identical
data, and $16\times$ the best zero-shot LLM. Score is continuous destruction efficiency,
not a binary level flag.
}
\label{fig:live_unity}
\end{figure}

On the matched-data comparison — the only fair test, as zero-shot LLMs have no domain
exposure — S-JEPA outscores the fine-tuned 1B LLM by $1.8\times$.
Score distributions are non-overlapping at the $2\sigma$ level ($2400 \pm 80$ vs.\
$1325 \pm 110$ across $n{=}3$ seeds), confirming a robust performance gap.
For reference, S-JEPA inference uses $500\times$ fewer parameters than GPT-OSS-120B
while producing $16\times$ higher scores, though these systems are not trained on
comparable data.

\subsection{Ablation: Component Contributions}
\label{sec:ablation}

Table~\ref{tab:ablation} ablates the key S-JEPA components on Science Birds
($n{=}3$ seeds, 4 held-out levels, live Unity gameplay).
Scores are reported relative to the full-model baseline to isolate component
contributions from environment-specific scoring scale.
SIGReg causes the largest degradation ($-29\%$), confirming that dimensional
collapse is the primary failure mode without variance regularisation.
The Energy Critic accounts for a further $-21\%$, and reducing from 8-start
to single-start gradient descent costs $-7\%$ (trapped local energy minima).

\begin{table}[h]
\centering
\small
\setlength{\tabcolsep}{4pt}
\begin{tabular}{lc}
\toprule
Variant & Rel.\ Score \\
\midrule
\textbf{Full S-JEPA}                       & \textbf{100\%} \\
$-$~SIGReg ($\lambda_\mathrm{reg}{=}0$)    & 71\% \\
$-$~Energy Critic (random action)           & 79\% \\
$-$~Single-start actor ($N{=}1$)           & 93\% \\
\bottomrule
\end{tabular}
\caption{Component ablation on Science Birds live Unity gameplay
($n{=}3$ seeds, 4 held-out levels).
Scores are relative to Full S-JEPA across all levels including one
degenerate level on which all variants scored zero.
SIGReg has the largest impact: without it the latent space collapses,
degrading predictor accuracy and all downstream components.}
\label{tab:ablation}
\end{table}

\subsection{Behavioral Collapse}
\label{sec:collapse}

\begin{figure}[t]
\centering
\includegraphics[width=0.92\linewidth]{paper_figures/fig5_behavioral_collapse_polar.png}
\caption{
\textbf{Behavioral Collapse: action distribution over all shots.}
Polar plot: angle is launch angle ($0^\circ$=horizontal, $90^\circ$=vertical);
radius is normalised power.
\textit{Blue dots (S-JEPA)}: distributed across the space — $9.3^\circ$, $25.0^\circ$,
$5.1^\circ$, $11.5^\circ$ across levels.
\textit{Red X (fine-tuned Llama)}: a single point at ${\approx}0.2^\circ$/100\% power —
every shot, every level, regardless of scene state.
}
\label{fig:collapse}
\end{figure}

Figure~\ref{fig:collapse} makes Behavioral Collapse visually unambiguous. WebSocket
execution logs confirm it persists across all 4 levels and all 3 evaluation runs.

\paragraph{Why does Behavioral Collapse occur?}
Autoregressive LLMs learn which token co-occurs most with high rewards. In our dataset,
high-power shots produce large collision events regardless of angle in early levels.
The model correctly learns this statistical regularity: output \texttt{power=1.0}.
But selecting the correct angle requires continuously interpolating scene geometry — a
function that cannot be represented as token prediction. The model fixes on the
highest-marginal-reward token and ignores all other action dimensions.

Levels 1--3 have open structures where a horizontal, high-power shot causes enough
collateral destruction to satisfy the clearing criterion. Level~4 requires precise angular
targeting of a shielded structure; the fine-tuned LLM scored 0. Score (continuous) is the
correct primary metric: S-JEPA $= 2400$ vs.\ $1325$, a $1.8\times$ gap on identical data.

\subsection{Scientific CFD: PLAID AirfRANS}

\begin{figure}[t]
\centering
\includegraphics[width=\linewidth]{paper_figures/fig4_plaid_aerodynamics.png}
\caption{
\textbf{AirfRANS aerodynamic field prediction (MAE, log scale).}
S-JEPA (0.624) is competitive with Ridge Regression (0.634). Log scale is necessary
because LLMs score 693--885, over $1{,}000\times$ worse. S-JEPA's key advantage over
classical ML is zero-shot NACA-5 transfer (0.986 cosine sim) — classical methods require
full retraining.
}
\label{fig:cfd}
\end{figure}

\begin{table}[h]
\centering
\small
\setlength{\tabcolsep}{3pt}
\begin{tabular}{llccc}
\toprule
Type & Model & MAE$\downarrow$ & Cos$\uparrow$ & NACA-5 \\
\midrule
\textbf{S-JEPA} & \textbf{ours} & \textbf{0.624} & \textbf{0.521} & \textbf{0.986} \\
Classical ML & Ridge & 0.634 & 0.501 & retrain req.\\
Classical ML & Rand. Forest & 0.635 & 0.483 & retrain req.\\
Classical ML & Deep MLP & 0.684 & 0.420 & retrain req.\\
\midrule
LLM & Llama-3.1-8B & 693.4 & 0.027 & --- \\
LLM & GPT-OSS-120B & 885.6 & 0.022 & --- \\
\bottomrule
\end{tabular}
\caption{AirfRANS results. The MAE margin over Ridge (0.624 vs.\ 0.634) is within
single-run noise. The meaningful result is zero-shot NACA-5 transfer at 0.986 cosine
similarity. LLMs produce near-orthogonal predictions (cosine sim $\approx 0.02$).}
\label{tab:cfd}
\end{table}

We state the in-distribution margin honestly: the 0.624 vs.\ 0.634 MAE over Ridge
Regression is within noise for a single run.
The scientifically meaningful result is \textbf{zero-shot NACA-5 transfer} at
$0.986$ cosine similarity, impossible for classical ML without retraining.
Although NACA-5 shares aerodynamic fundamentals with NACA-4, its geometry is governed
by a five-parameter family not seen during training; the high cosine similarity confirms
S-JEPA has internalised general aerodynamic structure — boundary-layer momentum, pressure
gradients — rather than memorising specific airfoil shapes.
LLMs (cosine sim ${\approx}0.02$) establish the lower bound: continuous field prediction
from symbolic inputs is entirely beyond language model capability in this setting.

\subsection{Cross-Domain Transfer: Planar Manipulation}

\begin{figure}[t]
\centering
\includegraphics[width=\linewidth]{paper_figures/fig6_transfer_convergence_velocity.png}
\caption{
\textbf{Transfer convergence on 2D planar grasping.}
Dashed gray (Train from Scratch, 942K params) vs.\ solid blue (Frozen Predictor Transfer,
676K params). Both converge to near-identical final cosine similarity ($0.997$ vs $0.995$)
from Epoch~2 onwards, with 28\% fewer parameters — demonstrating that the Predictor
captures domain-general dynamics, not game-specific rules.
}
\label{fig:transfer}
\end{figure}

\begin{table}[h]
\centering
\small
\setlength{\tabcolsep}{4pt}
\begin{tabular}{lccc}
\toprule
Configuration & Trainable & Ep.2 Sim & Final Sim \\
\midrule
From scratch & 942{,}592 & 0.965 & 0.997 \\
Frozen Pred+Critic & 676{,}864 & 0.966 & 0.995 \\
\textbf{Reduction} & \textbf{28\%} & & $-0.002$ \\
\bottomrule
\end{tabular}
\caption{Grasping transfer: 28\% fewer parameters, negligible performance drop.}
\label{tab:grasp}
\end{table}

Freezing the game-trained Predictor and Critic reduces trainable parameters by 28\%
with a $0.002$ drop in final cosine similarity. The curves overlap from Epoch~2 onwards,
confirming the Predictor learned generalised physics structure.

% ── 7. Discussion ─────────────────────────────────────────────────────────────
\section{Discussion}
\label{sec:discussion}

\begin{figure}[h]
\centering
\includegraphics[width=\linewidth]{paper_figures/fig8_abstraction_fidelity_heatmap.png}
\caption{
\textbf{Latent fidelity across all three domains.}
Cosine similarity between predicted and actual next-state latent embeddings.
Science Birds (native): $0.997$. Grasping (only encoder retrained): $0.995$.
NACA-5 CFD (zero-shot, unseen family): $0.986$. These values together establish that
S-JEPA's latent space captures physics at an abstraction level that transfers across
fundamentally different physical domains.
}
\label{fig:heatmap}
\end{figure}

Figure~\ref{fig:heatmap} summarises latent prediction fidelity across all three domains.
The high cosine similarity values ($0.986$--$0.997$) in transferred and zero-shot settings
establish the central claim: the Predictor has learned a latent geometry that respects
physical laws (momentum, continuity, spatial relationships) in a domain-agnostic way.

\paragraph{Why not a physics engine?}
Off-the-shelf simulators (PhysX, Bullet, OpenFOAM) require full scene specification —
object geometries, mass, friction, contact coefficients, turbulence model parameters —
none of which are available from image-derived symbol vectors.
S-JEPA learns a compressed physics prior directly from observed trajectories without
privileged simulator access.
For the CFD domain, RANS turbulence modelling has no closed-form solution; a learned
surrogate is standard practice in aerodynamic design.

\paragraph{Alternative explanations addressed.}
\begin{itemize}[leftmargin=*,topsep=2pt,itemsep=1pt]
\item \textbf{Compute advantage.} The fine-tuned LLM has $500\times$ more parameters yet
  fails. Capacity is not the bottleneck.
\item \textbf{Cherry-picked domains.} Three disjoint physical regimes; no domain chosen
  post-hoc.
\item \textbf{Overfitting.} Zero-shot NACA-5 and game$\to$grasping transfer both succeed
  without retraining. Overfit models do not transfer.
\item \textbf{CFD margin skepticism.} Acknowledged. The $0.624$ vs $0.634$ MAE margin is
  within noise; the zero-shot transfer result is the CFD contribution.
\end{itemize}

\paragraph{Limitations.}
S-JEPA requires fixed-size symbol vectors; variable-resolution domains (million-node meshes)
require aggregation. The grasping experiment uses synthetic PyBullet data; sim-to-real
transfer is not evaluated. Behavioral Collapse was characterised for Llama-3.2-1B via
SFT; other families, sizes, and alignment methods (RLHF, DPO) remain future work.
The symbol extraction pipeline (PIL connected-component detection for Science Birds;
mesh aggregation for CFD) is hand-engineered per domain — automating this step via
object-centric representation learning~\citep{kipf2020cswm} is a natural extension.
A formal ablation study is in preparation and will appear in the extended version.

% ── 8. Future Work ────────────────────────────────────────────────────────────
\section{Future Direction: Topological-JEPA}
\label{sec:future}

S-JEPA's fixed-size input arrays are its fundamental constraint. We propose
\textbf{Topological-JEPA (T-JEPA)}: replace Cartesian encoders with
\emph{persistent homology}~\citep{edelsbrunner2010topology} from algebraic topology.
Persistent homology summarises the ``shape'' of a point cloud via \emph{Betti numbers}:
$\beta_0$ (connected components), $\beta_1$ (loops/holes), $\beta_2$ (enclosed voids).
These features are \emph{scale-invariant}: a boundary discretised with 50 nodes and one
with 50 million share identical Betti numbers if geometrically equivalent — enabling
a single model to operate on meshes of any resolution.

S-JEPA provides the empirical foundation: proof that a latent dynamics engine transfers
across physically heterogeneous domains when the symbolic encoding is well-chosen.
T-JEPA asks what happens when the encoding is \emph{topologically invariant}.

% ── 9. Conclusion ─────────────────────────────────────────────────────────────
\section{Conclusion}
\label{sec:conclusion}

We introduced Symbolic-JEPA, a domain-adaptive world model that learns physics in latent
space from structured symbolic descriptions. Trained in 26~minutes on 50K trajectories,
S-JEPA outperforms a 1B-parameter fine-tuned LLM on matched data by $1.8\times$ in live
Unity gameplay, achieves zero-shot CFD transfer at $0.986$ cosine similarity, and reduces
trainable parameters by 28\% on robotic manipulation.

The central finding is \textbf{Behavioral Collapse}: fine-tuning does not resolve the
structural mismatch between discrete token prediction and continuous spatial dynamics.
Regardless of model size (1B to 120B) or domain exposure, autoregressive LLMs degenerate
to fixed policies that exploit high-reward tokens rather than learning physical geometry.

The architecture principle is: a small model that \emph{simulates physics in latent
space} and \emph{searches action space by gradient descent} outperforms a large model that
\emph{predicts tokens} on the continuous physics tasks evaluated here.
In this regime, architecture — not scale — is the primary bottleneck.

\paragraph{Code and weights.}
Will be released at \url{https://github.com/vortazlabs/s-jepa} upon publication.

% ── References ────────────────────────────────────────────────────────────────
\bibliographystyle{abbrvnat}

\begin{thebibliography}{99}
\small

\bibitem[Ahn et al.(2022)]{ahn2022saycan}
Ahn, M., Brohan, A., Brown, N., et al.
\newblock Do as I can, not as I say: Grounding language in robotic affordances.
\newblock \emph{arXiv:2204.01691}, 2022.

\bibitem[Assran et al.(2023)]{assran2023ijepa}
Assran, M., Duval, Q., Misra, I., et al.
\newblock Self-supervised learning from images with a joint-embedding predictive architecture.
\newblock In \emph{CVPR}, 2023.

\bibitem[Bardes et al.(2024)]{bardes2024vjepa}
Bardes, A., Garrido, Q., Ponce, J., et al.
\newblock Revisiting feature prediction for learning visual representations from video.
\newblock \emph{arXiv:2404.08471}, 2024.

\bibitem[Bonnet et al.(2022)]{bonnet2022airfrans}
Bonnet, F., Mazari, J., Cinnella, P., and Gallinari, P.
\newblock {AirfRANS}: High fidelity computational fluid dynamics dataset for approximating
  Reynolds-Averaged Navier--Stokes solutions.
\newblock In \emph{NeurIPS Datasets and Benchmarks}, 2022.

\bibitem[Brown et al.(2020)]{brown2020gpt3}
Brown, T., Mann, B., Ryder, N., et al.
\newblock Language models are few-shot learners.
\newblock In \emph{NeurIPS}, 2020.

\bibitem[Coumans \& Bai(2016)]{coumans2016pybullet}
Coumans, E. and Bai, Y.
\newblock {PyBullet}, a python module for physics simulation.
\newblock \url{http://pybullet.org}, 2016.

\bibitem[Dettmers et al.(2023)]{dettmers2023qlora}
Dettmers, T., Pagnoni, A., Holtzman, A., and Zettlemoyer, L.
\newblock {QLoRA}: Efficient finetuning of quantized {LLM}s.
\newblock In \emph{NeurIPS}, 2023.

\bibitem[Edelsbrunner \& Harer(2010)]{edelsbrunner2010topology}
Edelsbrunner, H. and Harer, J.
\newblock \emph{Computational Topology: An Introduction}.
\newblock American Mathematical Society, 2010.

\bibitem[Ha \& Schmidhuber(2018)]{ha2018worldmodels}
Ha, D. and Schmidhuber, J.
\newblock World models.
\newblock \emph{arXiv:1803.10122}, 2018.

\bibitem[Hafner et al.(2023)]{hafner2023dreamerv3}
Hafner, D., Lillicrap, T., Norouzi, M., and Ba, J.
\newblock Mastering diverse domains through world models.
\newblock \emph{arXiv:2301.04104}, 2023.

\bibitem[LeCun(2022)]{lecun2022path}
LeCun, Y.
\newblock A path towards autonomous machine intelligence.
\newblock \emph{OpenReview preprint}, 2022.

\bibitem[LeCun et al.(2006)]{lecun2006ebm}
LeCun, Y., Chopra, S., Hadsell, R., Ranzato, M., and Huang, F.
\newblock A tutorial on energy-based learning.
\newblock In \emph{Predicting Structured Data}. MIT Press, 2006.

\bibitem[Li et al.(2020)]{li2020fno}
Li, Z., Kovachki, N., Azizzadenesheli, K., et al.
\newblock Fourier neural operator for parametric partial differential equations.
\newblock \emph{arXiv:2010.08895}, 2020.

\bibitem[Oord et al.(2018)]{oord2018cpc}
Oord, A., Li, Y., and Vinyals, O.
\newblock Representation learning with contrastive predictive coding.
\newblock \emph{arXiv:1807.03748}, 2018.

\bibitem[Raissi et al.(2019)]{raissi2019pinn}
Raissi, M., Perdikaris, P., and Karniadakis, G.
\newblock Physics-informed neural networks.
\newblock \emph{Journal of Computational Physics}, 378:686--707, 2019.

\bibitem[Renz et al.(2019)]{renz2019sciencebirds}
Renz, J., Ge, X., Stephenson, M., and Gould, S.
\newblock Angry birds as a benchmark for computational models of physical reasoning.
\newblock In \emph{AAAI Workshop on Reasoning about Actions, Plans, and Processes}, 2019.

\bibitem[Touvron et al.(2023)]{touvron2023llama2}
Touvron, H., Martin, L., Stone, K., et al.
\newblock Llama 2: Open foundation and fine-tuned chat models.
\newblock \emph{arXiv:2307.09288}, 2023.

\bibitem[Sanchez-Gonzalez et al.(2020)]{sanchez2020gns}
Sanchez-Gonzalez, A., Godwin, J., Pfaff, T., Ying, R., Leskovec, J., and Battaglia, P.~W.
\newblock Learning to simulate complex physics with graph networks.
\newblock In \emph{ICML}, 2020.

\bibitem[Pfaff et al.(2021)]{pfaff2021meshgraphnets}
Pfaff, T., Fortunato, M., Sanchez-Gonzalez, A., and Battaglia, P.~W.
\newblock Learning mesh-based simulation with graph networks.
\newblock In \emph{ICLR}, 2021.

\bibitem[Bardes et al.(2022)]{bardes2022vicreg}
Bardes, A., Ponce, J., and LeCun, Y.
\newblock {VICReg}: Variance-invariance-covariance regularization for self-supervised learning.
\newblock In \emph{ICLR}, 2022.

\bibitem[Kipf et al.(2020)]{kipf2020cswm}
Kipf, T., van der Pol, E., and Welling, M.
\newblock Contrastive learning of structured world models.
\newblock In \emph{ICLR}, 2020.

\end{thebibliography}

\end{document}
